% GENERATED FILE. Edit canonical lessons, then run npm run generate:books.
% canonical-source-hash: fnv1a64:1d4fac3beca864eb
% canonical-lessons: ZH-C438-kaifa, ZH-C438-beibu, ZH-W438-bu-part, ZH-C438-xibu, ZH-C438-bufen, ZH-C438-quanbu

\chapter{The North}
\label{ch:zh-the-north}
\hlchaptermodality{438}

\begin{chapteropening}
\textbf{By the end of this chapter:} \emph{I can talk about developing an area, the northern and western parts, a part of the visitors and all of the luggage.}

Complete the last lesson of chapter 438: I can talk about developing an area, the northern and western parts, a part of the visitors and all of the luggage.
\end{chapteropening}

\section[\texorpdfstring{\zh{开发} (kāifā)}{开发 (kāifā)}]{\zh{开发} (kāifā) — to open up land for use}

\label{lesson:ZH-C438-kaifa}



\begin{quote}
Before the new one: say the Mandarin for a high plateau, then the Mandarin for to emigrate.
\end{quote}

\subsection*{You'll want to know: \zh{开发}}
\textbf{\zh{开发}} — \emph{kāifā} — \textquotedblleft{}to open up land for use\textquotedblright{}.

\begin{quote}
\textbf{\zh{这个山区还没有开发}} \emph{zhège shānqū hái méiyǒu kāifā} — This mountain area hasn't been developed yet.
\end{quote}

\begin{cousinweb}
\textbf{\zh{开发}} (\emph{kāifā}) is built from \zh{开} 'to open' and \zh{发} 'to set going', both inherited from Middle Chinese, \emph{khoj} and \emph{pjot} (Guangyun, in Baxter's transcription), which glosses the second 'to start' and 'to raise'.
\end{cousinweb}

\begin{grammarlens}[title={What you've built}]
One more word for getting around.
\end{grammarlens}

\subsection*{Your turn}
\begin{itemize}
\raggedright
  \item \emph{Say it:} \emph{kāifā}
  \item \emph{Say it:} \emph{kāifā}, once more
  \item \emph{Say it:} \emph{zhège shānqū hái méiyǒu kāifā} — This mountain area hasn't been developed yet.
  \item \emph{Recall:} say the Mandarin for a high plateau, then the Mandarin for to emigrate, then say \emph{kāifā} again
\end{itemize}

\begin{tcolorbox}[breakable,colback=teal!4,colframe=teal!35!black,title={Before you move on}]
What does \zh{开发} mean? (\textquotedblleft{}To open up land for use\textquotedblright{}.) Say it once more.
\end{tcolorbox}
\section[\texorpdfstring{\zh{北部} (běibù)}{北部 (běibù)}]{\zh{北部} (běibù) — northern part}

\label{lesson:ZH-C438-beibu}



\begin{quote}
Before the new one: say the Mandarin for to emigrate, then the Mandarin for to open up land for use.
\end{quote}

\subsection*{You'll want to know: \zh{北部}}
\textbf{\zh{北部}} — \emph{běibù} — \textquotedblleft{}northern part\textquotedblright{}.

\begin{quote}
\textbf{\zh{我们坐火车去中国北部}} \emph{wǒmen zuò huǒchē qù Zhōngguó běibù} — We're taking the train up to the north of China.
\end{quote}

The new character, in pieces: \zh{部} — ten strokes in two pieces: on the left five strokes, a dot, a bar, two short strokes and a long bar, over \zh{口} kǒu; on the right the right piece of \zh{那}, two strokes. The next lesson puts it in your hand.

\begin{cousinweb}
\textbf{\zh{北部}} (\emph{běibù}) is built from \zh{北} 'north' and \zh{部} 'a part', both inherited from Middle Chinese, \emph{pok} and \emph{buX} (Guangyun, in Baxter's transcription), which glosses the second as a body of troops; the sense 'a part of a whole' is not given there.
\end{cousinweb}

\begin{grammarlens}[title={What you've built}]
One more word for getting around.
\end{grammarlens}

\subsection*{Your turn}
\begin{itemize}
\raggedright
  \item \emph{Say it:} \emph{běibù}
  \item \emph{Say it:} \emph{běibù}, once more
  \item \emph{Say it:} \emph{wǒmen zuò huǒchē qù Zhōngguó běibù} — We're taking the train up to the north of China.
  \item \emph{Recall:} say the Mandarin for to emigrate, then the Mandarin for to open up land for use, then say \emph{běibù} again
\end{itemize}

\begin{tcolorbox}[breakable,colback=teal!4,colframe=teal!35!black,title={Before you move on}]
What does \zh{北部} mean? (\textquotedblleft{}Northern part\textquotedblright{}.) Say it once more.
\end{tcolorbox}
\section[\texorpdfstring{\zh{部} (bù)}{部 (bù)}]{\zh{部} — five strokes over \zh{口}, then the right piece of \zh{那}}

\label{lesson:ZH-W438-bu-part}



\begin{quote}
The word first, then the character inside it.

\begin{itemize}
\raggedright
  \item \emph{Recall:} say \emph{běibù} — \textbf{R1}, one lesson back, before it decays
\end{itemize}

\begin{quote}
\zh{那} \zh{口}
\end{quote}

\emph{Read it:} each one aloud, with nothing said first — \emph{nà}, \emph{kǒu}
\end{quote}

\subsection*{Script — five strokes over \zh{口}, then the right piece of \zh{那}}
\hlblockfigure{\detokenize{figures/ZH-W438-bu-part-filmstrip.pdf}}{How \zh{部} is written, stroke by stroke}

\begin{quote}
\zh{部}
\end{quote}

\textbf{Two pieces, left before right.} Count the pieces, not the strokes.

\textbf{The left piece, eight strokes, in two parts.} On top, five strokes: a dot, a bar, then a short dot on the left and a short stroke falling to the left on the right, and a long bar under them. Under them, \zh{口} kǒu, three strokes.

\textbf{The right piece, two strokes}, as on the right of \zh{那}: a stroke that runs right, bends back and curls down to a hook, then a long vertical.

5 + 3 + 2 = ten.

\subsection*{Your turn}
\begin{itemize}
\raggedright
  \item \emph{Write it:} \zh{部}, all ten, with the model in front of you
  \item \emph{Cover:} write it again from memory, and count the pieces, not the strokes
\end{itemize}

\begin{tcolorbox}[breakable,colback=teal!4,colframe=teal!35!black,title={Before you move on}]
How many strokes in \zh{部}? (\textbf{Ten}.) What sits under the five strokes on top? (\textbf{\zh{口}}.) Where has your hand made the right piece before? (\textbf{On the right of \zh{那}}.)

Source: \href{https://raw.githubusercontent.com/chanind/hanzi-writer-data/68d10a4b21150cae5e1ebbd223eed289cf32d90c/data/\%E9\%83\%A8.json}{Hanzi Writer Data \zh{部}, pinned stroke paths and medians}.
\end{tcolorbox}
\section[\texorpdfstring{\zh{西部} (xībù)}{西部 (xībù)}]{\zh{西部} (xībù) — western part; the western regions}

\label{lesson:ZH-C438-xibu}



\begin{quote}
Before the new one: say the Mandarin for to open up land for use, then the Mandarin for the northern part.
\end{quote}

\subsection*{You'll want to know: \zh{西部}}
\textbf{\zh{西部}} — \emph{xībù} — \textquotedblleft{}western part; the western regions\textquotedblright{}.

\begin{quote}
\textbf{\zh{中国西部的山很高}} \emph{Zhōngguó xībù de shān hěn gāo} — The mountains in western China are very high.
\end{quote}

\begin{cousinweb}
\textbf{\zh{西部}} (\emph{xībù}) is built from \zh{西} 'west' and \zh{部} 'a part', both inherited from Middle Chinese, \emph{sej} and \emph{buX} (Guangyun, in Baxter's transcription).
\end{cousinweb}

\begin{grammarlens}[title={What you've built}]
One more word for getting around.
\end{grammarlens}

\subsection*{Your turn}
\begin{itemize}
\raggedright
  \item \emph{Say it:} \emph{xībù}
  \item \emph{Say it:} \emph{xībù}, once more
  \item \emph{Say it:} \emph{Zhōngguó xībù de shān hěn gāo} — The mountains in western China are very high.
  \item \emph{Recall:} say the Mandarin for to open up land for use, then the Mandarin for the northern part, then say \emph{xībù} again
\end{itemize}

\begin{tcolorbox}[breakable,colback=teal!4,colframe=teal!35!black,title={Before you move on}]
What does \zh{西部} mean? (\textquotedblleft{}Western part; the western regions\textquotedblright{}.) Say it once more.
\end{tcolorbox}
\section[bùfen]{bùfen — part; section (of a whole)}

\label{lesson:ZH-C438-bufen}



\begin{quote}
Before the new one: say the Mandarin for the northern part, then the Mandarin for the western part.
\end{quote}

\subsection*{You'll want to know: bùfen}
\emph{bùfen} — \textquotedblleft{}part; section (of a whole)\textquotedblright{}.

\begin{quote}
\emph{dà bùfen yóukè zhù zài běibù} — Most of the visitors are staying in the north.
\end{quote}

Its characters come later; for now it is a word for the ear.

\begin{cousinweb}
\emph{bùfen} is built from \zh{部} 'a section' and \emph{fèn} 'a share', both inherited from Middle Chinese, \emph{buX} and \emph{bjunH} (Guangyun, in Baxter's transcription), which glosses the second as a portion that is measured out; in this word it is said light.
\end{cousinweb}

\begin{grammarlens}[title={What you've built}]
One more word for getting around.
\end{grammarlens}

\subsection*{Your turn}
\begin{itemize}
\raggedright
  \item \emph{Say it:} \emph{bùfen}
  \item \emph{Say it:} \emph{bùfen}, once more
  \item \emph{Say it:} \emph{dà bùfen yóukè zhù zài běibù} — Most of the visitors are staying in the north.
  \item \emph{Recall:} say the Mandarin for the northern part, then the Mandarin for the western part, then say \emph{bùfen} again
\end{itemize}

\begin{tcolorbox}[breakable,colback=teal!4,colframe=teal!35!black,title={Before you move on}]
What does bùfen mean? (\textquotedblleft{}Part; section (of a whole)\textquotedblright{}.) Say it once more.
\end{tcolorbox}
\section[\texorpdfstring{\zh{全部} (quánbù)}{全部 (quánbù)}]{\zh{全部} (quánbù) — all of it; the whole}

\label{lesson:ZH-C438-quanbu}



\begin{quote}
Before the new one: say the Mandarin for the western part, then the Mandarin for a part of a whole.
\end{quote}

\subsection*{You'll want to know: \zh{全部}}
\textbf{\zh{全部}} — \emph{quánbù} — \textquotedblleft{}all of it; the whole\textquotedblright{}.

\begin{quote}
\emph{wǒmen de xíngli quánbù fàng zài chē lǐ} — All our luggage is in the car.
\end{quote}

\begin{cousinweb}
\textbf{\zh{全部}} (\emph{quánbù}) is built from \zh{全} 'whole' and \zh{部} 'a part', both inherited from Middle Chinese, \emph{dzjwen} and \emph{buX} (Guangyun, in Baxter's transcription), which glosses the first 'complete'.
\end{cousinweb}

\begin{grammarlens}[title={What you've built}]
One more word for getting around.
\end{grammarlens}

\subsection*{Your turn}
\begin{itemize}
\raggedright
  \item \emph{Say it:} \emph{quánbù}
  \item \emph{Say it:} \emph{quánbù}, once more
  \item \emph{Say it:} \emph{wǒmen de xíngli quánbù fàng zài chē lǐ} — All our luggage is in the car.
  \item \emph{Recall:} say the Mandarin for the western part, then the Mandarin for a part of a whole, then say \emph{quánbù} again
\end{itemize}

\begin{tcolorbox}[breakable,colback=teal!4,colframe=teal!35!black,title={Before you move on}]
What does \zh{全部} mean? (\textquotedblleft{}All of it; the whole\textquotedblright{}.) Say it once more.
\end{tcolorbox}
